\chapter{Sobolev spaces on surfaces}

\cite[p. 299]{dziuk_finite_2013} Let $\surflisse \in \Cont^2$ for the following. For $p \in \intervalleff{1}{\infty}$ we let $\L^p(\surflisse)$ denote the space of functions $\fonctionens[f]{\surflisse}{\R}$ which are measurable with respect to the surface measure $\d A$ ($\d A$ in connection with an integral over $\surflisse$ denotes the $n$-dimensional Hausdorff measure) and have finite norm, where the norm is defined by 
\[
\norme{f}_{\L^p(\surflisse)} \defeq \bigg( \int_{\surflisse} \abs{f}^p \d A \bigg)^{1/p}
\]
for $p < \infty$, and for $p = \infty$ we mean the essential supremum norm. The space $\L^p(\surflisse)$ is a Banach space and for $p = 2$ a Hilbert space. For $1 \leqslant p < \infty$ the space $\Cont^0(\surflisse)$ dans $\Cont^1(\surflisse)$ are dense in $\L^p(\surflisse)$. 

\begin{theo}[Poincaré’s inequality {\cite[Sec. 2.4]{dziuk_finite_2013}}]
    Assume that $\surflisse \in \Cont^3$ and $1 \leqslant p < \infty$. Then there is a constant $c$ such that, for every function $f \in \H^{1,p}(\surflisse)$ with $\int_{\surflisse} f \d A = 0$, we have the inequality
    \[
    \norme{f}_{\L^p(\surflisse)} \leqslant c \norme{\grad_{\surflisse} f}_{\L^p(\surflisse)}.
    \]
\end{theo}